\section{傅里叶级数的渐进特性}\label{sec:Asymptotic_Behaviour}

\subsection{狄利克雷核}\label{sub:dirichlet_kernel}
本小节引入狄利克雷核并介绍它的一些性质，并在下一小节用它来辅助我们对傅里叶级数收敛性的讨论。直接由狄利克雷条件证明收敛性是很困难的，并且需要更加专业的分析学工具，我们将给出更强的条件下的证明，并从这个更强的条件获得一些额外的性质。

研究傅里叶级数的渐进特性时，一个非常好用的工具是\textbf{狄利克雷核}(Dirichlet kernel)：
\begin{definition}[狄利克雷核]
    \begin{equation}
        D_N(t)=\sum_{k=-N}^{N}\e^{\im k\omega t}=1+\sum_{k=1}^{N}\paren{\e^{\im k\omega t}+\e^{-\im k\omega t}}=1+2\sum_{k=1}^{N}\cos(k\omega t),
    \end{equation}
    它是依赖于所研究函数的周期$T$的，但简便起见，在符号$D_N(t)$中不体现这一点。
\end{definition}

我们可以用等比数列求和或积化和差裂项的方法化简$D_N(t)$：
\begin{align}
    D_N(t)&=\sum_{k=-N}^{N}\e^{\im k\omega t} \nonumber \\
    &=\e^{-\im N\omega t}\frac{1-\e^{\im(2N+1)\omega t}}{1-\e^{\im\omega t}} \nonumber \\
    &=\frac{\e^{\im(N+1)\omega t}-\e^{-\im N\omega t}}{\e^{\im\omega t}-1}, \\
    D_N(t)&=1+\sum_{k=1}^{N}\paren{\e^{\im k\omega t}+\e^{-\im k\omega t}} \nonumber \\
    &=1+2\sum_{k=1}^{N}\cos(k\omega t) \nonumber \\
    &=1+\frac{2}{\sin\paren{\frac{\omega t}{2}}}\sum_{k=1}^{N}\cos(k\omega t)\sin\paren{\frac{\omega t}{2}} \nonumber \\
    &=1+\frac{1}{\sin\paren{\frac{\omega t}{2}}}\sum_{k=1}^{N}\sbr{\sin\paren{\paren{k+\frac{1}{2}}\omega t}-\sin\paren{\paren{k-\frac{1}{2}}\omega t}} \nonumber \\
    &=1+\frac{\sin\paren{\paren{N+\frac{1}{2}}\omega t}-\sin\paren{\frac{\omega t}{2}}}{\sin\paren{\frac{\omega t}{2}}} \nonumber \\
    &=\frac{\sin\paren{\paren{N+\frac{1}{2}}\omega t}}{\sin\paren{\frac{\omega t}{2}}}.
\end{align}
这两种结果是相符的，读者可自行验证，并且可以从后一结果想象出狄利克雷核的函数图像，它被$\pm 1/\sin\paren{\frac{\omega t}{2}}$包络并高速振荡。函数图像如下。
\begin{figure}[H]
    \centering
    \includegraphics[width=0.8\textwidth]{Figure_3}
    \caption{狄利克雷核$D_N(t)$的函数图像}
\end{figure}

傅里叶级数的部分和是狄利克雷核与被展开函数的周期函数卷积（\cref{def:period_conv}）：
\begin{proposition}
    \begin{equation}
        S_N^f(t)=\frac{1}{T}\int_{T}f(\tau)D_N(t-\tau)\,\dd\tau=\frac{1}{T}f*D_N(t).
    \end{equation}
\end{proposition}
\begin{proof}
    \begin{align}
        S_N^f(t)&=\sum_{-N}^{N}c_k \e^{\im k\omega t} \nonumber \\
        &=\sum_{-N}^{N}\paren{\frac{1}{T}\int_{T}f(\tau)\e^{-\im k\omega \tau}\,\dd\tau}\e^{\im k\omega t} \nonumber \\
        &=\frac{1}{T}\int_{T}f(\tau)\sum_{k=-N}^{N}\e^{\im k\omega (t-\tau)}\,\dd\tau \nonumber \\
        &=\frac{1}{T}\int_{T}f(\tau)D_N(t-\tau)\,\dd\tau \nonumber \\
        &=\frac{1}{T}\int_{T}f(t-\tau)D_N(\tau)\,\dd\tau \qquad (\tau\to t-\tau) \nonumber \\
        &=\frac{1}{T}\int_{T}f(t+\tau)D_N(\tau)\,\dd\tau. \qquad (\tau\to t+\tau)
    \end{align}
\end{proof}

在讨论傅里叶级数的收敛性前，先给出两个引理。第一个引理表明狄利克雷核在半周期上积分值为$\displaystyle\frac{T}{2}$，在证明傅里叶级数的逐点收敛性时将用到它。
\begin{lemma}
    \begin{equation}
        \int_{-\frac{T}{2}}^{0}D_N(t)\,\dd t=\int_{0}^{\frac{T}{2}}D_N(t)\,\dd t=\frac{T}{2}.
    \end{equation}
\end{lemma}
\begin{proof}
    \begin{align}
        D_N(t)&=1+2\sum_{k=1}^{N}\cos(k\omega t), \\
        \int_{0}^{\frac{T}{2}}D_N(t)\,\dd t&=\int_{0}^{\frac{T}{2}}\paren{1+\sum_{k=1}^{N}\cos(k\omega t)}\,\dd t \nonumber \\
        &=\frac{T}{2}+\sum_{k=1}^{N}\evalat{\frac{\sin(k\omega t)}{k\omega}}{0}{\frac{T}{2}} \nonumber \\
        &=\frac{T}{2}+\frac{1}{\omega}\sum_{k=1}^{N}\frac{\sin(k\pi)}{k}=\frac{T}{2}.
    \end{align}
    $D_N(t)$是偶函数，得证。
\end{proof}

\begin{lemma}[贝塞尔不等式]\label{lemma:bessel_inequality}
    设$f\in L^2([0,T]),\ c_n=\displaystyle\frac{1}{T}\int_{T}f(t)\e^{-\im n\omega t}\,\dd t$，则
    \begin{equation}
        \sum_{n=-\infty}^{\infty}|c_n|^2\leq \frac{1}{T}\int_{T}|f(t)|^2\,\dd t=\frac{1}{T}\norm{f}_2^2.
    \end{equation}
\end{lemma}
它给出了傅里叶系数平方和的上界的估计。由此可以看出$c_n\to 0,\ n\to\infty$。
\begin{proof}
    \begin{align}
        \abs{f(t)-\sum_{n=-N}^{N}c_n \e^{\im n\omega t}}^2&=\paren{f(t)-\sum_{n=-N}^{N}c_n \e^{\im n\omega t}}\overline{\paren{f(t)-\sum_{n=-N}^{N}c_n \e^{\im n\omega t}}} \nonumber \\
        &=\paren{f(t)-\sum_{n=-N}^{N}c_n \e^{\im n\omega t}}\paren{\overline{f(t)}-\sum_{n=-N}^{N}\overline{c_n}\e^{-\im n\omega t}} \nonumber \\
        &=|f(t)|^2-\sum_{n=-N}^{N}\paren{\overline{c_n}f(t)\e^{\im n\omega t}+c_n\overline{f(t)}\e^{-\im n\omega t}}+\sum_{m,n=-N}^{N}c_m\overline{c_n}\e^{\im(m-n)\omega t}.
    \end{align}
    将上式在一个周期上积分，我们知道
    \begin{equation}
        \int_{T}f(t)\e^{\im n\omega t}\,\dd t=Tc_n,\quad \int_{T}\e^{\im(m-n)\omega t}\,\dd t=\begin{dcases}
            0, & \text{if } m\neq n, \\
            T, & \text{if } m=n,
        \end{dcases}
    \end{equation}
    故
    \begin{align}
        &\int_{T}|f(t)|^2\,\dd t-\sum_{n=-N}^{N}\paren{\overline{c_n}\int_{T}f(t)\e^{\im n\omega t}\,\dd t+c_n\int_{T}\overline{f(t)}\e^{-\im n\omega t}\,\dd t} \nonumber \\
        &\quad+\sum_{m,n=-N}^{N}c_m\overline{c_n}\int_{T}\e^{\im(m-n)\omega t}\,\dd t \nonumber \\
        &=\int_{T}|f(t)|^2\,\dd t-T\sum_{n=-N}^{N}(\overline{c_n} c_n+c_n\overline{c_n})+T\sum_{n=-N}^{N}\overline{c_n}c_n \nonumber \\
        &=\int_{T}|f(t)|^2\,\dd t-T\sum_{n=-N}^{N}|c_n|^2.
    \end{align}
    这是非负函数的积分，积分值非负，即
    \begin{equation}
        \sum_{-\infty}^{\infty}|c_n|^2\leq \frac{1}{T}\int_{T}|f(t)|^2\,\dd t<\infty.
    \end{equation}
\end{proof}

\subsection{分段光滑条件下的一致收敛性}\label{sub:uniform_converge}
直接由狄利克雷条件证明逐点收敛性需要很专业的分析学工具，但我们可以适当地加强狄利克雷条件，以得到更直观、简单的结论\cite{folland}。分段光滑的定义，见\cref{def:piecewice_smooth}。

我们研究的多数函数是分段光滑的。分段光滑会带来一些额外的性质，例如，可以相对简单地证明逐点收敛性：
\begin{theorem}[傅里叶级数的逐点收敛性]
    设$f\in PS(\mathbb{R})$是周期函数，其傅里叶级数的部分和为$S_N^f(t)$，则
    \begin{equation}
        \lim_{N\to\infty}S_N^f(t_0)=\frac{f(t_0+)+f(t_0-)}{2}.
    \end{equation}
\end{theorem}
\begin{proof}
    \begin{align}
        S_N^f(t_0)-\frac{f(t_0^+)+f(t_0^-)}{2}&=\frac{1}{T}\left(\int_{T}f(t_0-\tau)D_N(\tau)\,\dd\tau\right. \nonumber \\
        &\qquad-\left.\int_{0}^{\frac{T}{2}}f(t_0^+)D_N(\tau)\,\dd\tau-\int_{-\frac{T}{2}}^{0}f(t_0^-)D_N(\tau)\,\dd\tau\right) \nonumber \\
        &=\frac{1}{T}\left(\int_{0}^{\frac{T}{2}}(f(t_0-\tau)-f(t_0^+))D_N(\tau)\,\dd\tau\right. \nonumber \\
        &\qquad+\left.\int_{-\frac{T}{2}}^{0}(f(t_0-\tau)-f(t_0^-))D_N(\tau)\,\dd\tau\right),
    \end{align}
    \begin{equation}
        S_N^f(t_0)-\frac{f(t_0^+)+f(t_0^-)}{2}=\frac{1}{T}\int_{T}g(t)\paren{\e^{\im(N+1)\omega t}-\e^{-\im N\omega t}}\,\dd t,
    \end{equation}
    其中
    \begin{equation}
        g(t):=\begin{dcases}
            \frac{f(t_0+t)-f(t_0^-)}{\e^{\im\omega t}-1}, & \text{if } -\frac{T}{2}<t<0, \\
            \frac{f(t_0+t)-f(t_0^+)}{\e^{\im\omega t}-1}, & \text{if } 0<t<\frac{T}{2}.
        \end{dcases}
    \end{equation}
    由洛必达法则，$t\to 0$时，
    \begin{equation}
        \lim_{t\to 0^+}g(t)=\lim_{t\to 0^+}\frac{f(t_0+t)-f(t_0^+)}{\e^{\im\omega t}-1}=\lim_{t\to 0^+}\frac{f'(t_0+t)}{\im\omega \e^{\im\omega t}}=\frac{f'(t_0^+)}{\im\omega},
    \end{equation}
    $t\to 0^-$时同理。故$g$分段连续，当然是平方可积的，由贝塞尔不等式，$g(t)$的傅里叶系数平方和收敛，通项趋于$0$，$S_N^f(t_0)-(f(t_0^+)+f(t_0^-))/2=C_{-(N+1)}-C_N\to 0$，得证。
\end{proof}

在分段光滑的条件下，容易得到$f'(t)$的傅里叶系数，注意微积分基本定理可以分区间使用：
\begin{proposition}[导函数的傅里叶系数]
    \label{prop:derivative_fourier_series}
    设周期函数$f(t)\in PS(\mathbb{R})$的基波角频率为$\omega$，三角形式的傅里叶系数为$\{a_n\}_{n=0}^\infty$和$\{b_n\}_{n=1}^\infty$，复指数形式的傅里叶系数为$\{c_n\}_{n\in\mathbb{Z}}$，则其导函数$f'(t)$的对应的傅里叶系数为：
    \begin{equation}
        a'_n=n\omega b_n,\quad b'_n=-n\omega a_n,\quad c'_n=\im n\omega c_n.
    \end{equation}
\end{proposition}
\begin{proof}
    以$c_n$为例：
    \begin{align}
        c'_n&=\frac{1}{T}\int_{T}f'(t)\e^{-\im n\omega t}\,\dd t \nonumber \\
        &=\frac{1}{T}\evalat{f(t)\e^{-\im n\omega t}}{0}{T}+\im n\omega\int_{T}f(t)\e^{-\im n\omega t}\,\dd t \nonumber \\
        &=\im n\omega c_n.
    \end{align}
\end{proof}

$f$的原函数$F$的傅里叶系数同理，并且只要\textbf{分段连续}（\cref{def:piecewise_continuos}）即可保证$f$可积，但是我们必须保证$F$是周期函数，这要求$f$的直流分量为$0$：
\begin{equation}
    F(t+T)-F(t)=\int_{T}f(t)\,\dd t=Tc_0=0,\quad c_0=0.
\end{equation}
此时，用刚刚得到的\cref{prop:derivative_fourier_series}就可直接得到$F$的傅里叶系数：
\begin{proposition}[原函数的傅里叶系数]
    \label{prop:integral_fourier_series}
    设周期函数$f(t)\in PS(\mathbb{R})$的基波角频率为$\omega$，三角形式的傅里叶系数为$\{a_n\}_{n=0}^\infty$和$\{b_n\}_{n=1}^\infty$，复指数形式的傅里叶系数为$\{c_n\}_{n\in\mathbb{Z}}$，且直流分量$a_0/2=c_0=0$，则其原函数$F$的对应的傅里叶系数为：
    \begin{equation}
        A_n=\frac{a_n}{n\omega},\quad B_n=\frac{b_n}{n\omega},\quad C_n=\frac{c_n}{\im n\omega}.
    \end{equation}
    直流分量$A_0/2=C_0$仅依赖于不定积分的常数$C$。
\end{proposition}

分段光滑还能推出傅里叶级数的\textbf{一致收敛性}：
\begin{theorem}\label{thm:uniform_convergency}
    \begin{equation}
        f\in PS([0,T])\implies S_N^f(t)\rightrightarrows f(t).
    \end{equation}
\end{theorem}

%\begin{theorem}[M判别法]
%    对于函数项级数$\sum_{n=1}^{\infty}f_n(x)$，如果存在正项级数$\sum_{n=1}^{\infty}M_n<\infty$使得在区间E上$|f_n(x)|<M_n$，则$\sum_{n=1}^{\infty}f_n(x)$在E上绝对收敛且一致收敛。
%\end{theorem}
\begin{proof}
    使用魏尔斯特拉斯M判别法，要证明一致收敛性，只需证明
    \begin{equation}
        \sum_{n=-\infty}^{\infty}|c_n \e^{\im n\omega t}|=\sum_{n=-\infty}^{\infty}|c_n|<\infty.
    \end{equation}
    记$c'_n$为$f'$的傅里叶系数，$c'_n=\im n\omega c_n$，
    \begin{align}
        \sum_{n=-\infty}^{\infty}|c_n|&=|c_0|+\sum_{n\neq 0}\abs{\frac{c'_n}{n\omega}} \nonumber \\
        &\leq |c_0|+\frac{1}{\omega}\paren{\sum_{n\neq 0}\frac{1}{n^2}}^{1/2}\paren{\sum_{n\neq 0}|c'_n|^2}^{1/2}<\infty.
    \end{align}
    最后一步使用了柯西-施瓦兹不等式。
\end{proof}

请读者思考：我们探究了指数形式傅里叶级数收敛的条件，对于三角函数形式的傅里叶级数应该怎么办？

\subsection{收敛速度的估计}\label{sub:converge_speed}

\begin{theorem}\label{thm:converge_speed}
    设$f\in C^p(\mathbb{R})\ (p\geq 1)$是周期函数，其傅里叶系数为$c_k$，部分和
    \begin{equation}
        S_N^f(t)=\sum_{k=-N}^{N}c_k \e^{\im k\omega t}
    \end{equation}
    在$\mathbb{R}$上逐点收敛、内闭一致收敛，且
    \begin{equation}
        \max\abs{f(t)-S_N^f(t)}<\frac{C}{N^{p-0.5}},
    \end{equation}
    其中$C$为常数。
\end{theorem}
\begin{proof}
    内闭一致收敛性是\cref{thm:uniform_convergency}的推论。根据\cref{prop:derivative_fourier_series}，$f^{(p)}(t)$的复指数形式的傅里叶系数为
    \begin{equation}
        c^{(p)}_k=(\im k\omega)^p c_k,
    \end{equation}
    因此
    \begin{align}
        \max\abs{f(t)-S_N^f(t)}&=\max\abs{\sum_{k=N+1}^{\infty}c_k\e^{\im k\omega t}+\sum_{k=-\infty}^{-(N+1)}c_k\e^{\im k\omega t}} \nonumber \\
        &\leq\sum_{|k|>N}|c_k|=\sum_{|k|>N}\frac{|c_k^{(p)}|}{(k\omega)^p}.
    \end{align}
    由柯西不等式可得：
    \begin{align}
        \sum_{|k|>N}\frac{|c_k^{(p)}|}{(k\omega)^p}&\leq\paren{\sum_{|k|>N}\frac{1}{(k\omega)^{2p}}}^{1/2}\paren{\sum_{|k|>N}|c^{(p)}_k|^2}^{1/2} \nonumber \\
        &<\frac{1}{\omega^{p}}\paren{\sum_{|k|>N}\frac{1}{k^{2p}}}^{1/2}\paren{\sum_{k\in\mathbb{Z}}|c^{(p)}_k|^2}^{1/2}.
    \end{align}
    $f\in C^p(\mathbb{R})$表明$f^{(p)}(t)$是连续函数，因此$f^{(p)}(t)\in L^2([0,T])$，由\cref{lemma:bessel_inequality}可得：
    \begin{equation}
        \sum_{k\in\mathbb{Z}}\abs{c^{(p)}_k}^2\leq\frac{1}{T}\norm{f^{(p)}}_2^2<\infty,
    \end{equation}
    因此
    \begin{equation}
        \sum_{|k|>N}\frac{|c_k^{(p)}|}{(k\omega)^p}<\frac{\norm{f^{(p)}}_2}{\sqrt{T}\omega^{p}}\paren{\sum_{|k|>N}\frac{1}{k^{2p}}}^{1/2}.
    \end{equation}
    利用
    \begin{equation}
        \int_{k}^{k+1}\frac{\dd x}{x^{2p}}<\frac{1}{k^{2p}}<\int_{k-1}^{k}\frac{\dd x}{x^{2p}}
    \end{equation}
    对$k$求和，得到
    \begin{equation}
        \sum_{k=N+1}^{\infty}\frac{1}{k^{2p}}<\int_{N}^{\infty}\frac{\dd x}{x^{2p}}=\evalat{-\frac{1}{(2p-1)x^{2p-1}}}{N}{\infty}=\frac{1}{(2p-1)N^{2p-1}},
    \end{equation}
    因此
    \begin{equation}
        \max\abs{f(t)-S_N^f(t)}<\frac{2\norm{f^{(p)}}_2}{\sqrt{(2p-1)T}\omega^{p}}\cdot\frac{1}{N^{p-0.5}}.
    \end{equation}
\end{proof}
